\section{Introduction}
\label{sec:introduction}

Machine learning is increasingly used to support decisions in domains such as logistics, energy, healthcare, finance, and public-sector planning \citep{placeholder_ml_decision_making}. This recent progress can make automated decision making seem like a new phenomenon. In fact, computers have been used to make high-stakes operational decisions for decades. Since the development of linear programming and the simplex method, classically associated with Dantzig \citep{placeholder_dantzig_linear_programming}, mathematical programming has provided a rigorous framework for computing decisions by solving large constrained optimization problems.

Today, mathematical programming is central to the operation of many critical systems. Linear, integer, stochastic, and conic optimization models are used in electric power systems, supply chains, transportation, revenue management, finance, emergency response, and many other areas \citep{placeholder_energy_optimization,placeholder_supply_chain_optimization,placeholder_transportation_optimization,placeholder_finance_optimization,placeholder_healthcare_optimization}. These models are attractive because they explicitly encode feasibility, resource limits, network structure, and other domain constraints. At the same time, these same constraints make many mathematical programming problems difficult to address with standard deep learning methods. Neural networks can be excellent predictors, but enforcing hard feasibility, integrality, and recourse structure remains challenging \citep{placeholder_ml_constraints,placeholder_learning_constrained_decisions}.

This has motivated a growing line of work at the interface of machine learning and mathematical programming. Rather than replacing optimization models with neural networks, hybrid methods use learning to improve the inputs, approximations, or solution procedures of optimization problems \citep{placeholder_ml_for_optimization_survey}. Such methods aim to combine the statistical flexibility of machine learning with the feasibility guarantees and modeling structure of mathematical programming. Recent work has shown that this combination can lead to substantial performance improvements on large practical problems \citep{placeholder_bouvier_real_life_optimization,placeholder_hybrid_optimization_applications}. Given the societal importance of mathematical programming, even modest improvements in these systems can have significant practical impact.

A particularly important setting is stochastic programming. In a two-stage stochastic program, a decision maker first chooses an action before uncertainty is realized. After the uncertainty is observed, a second-stage recourse decision is made to correct or adapt the original decision. For example, an energy operator may commit generation capacity before demand and renewable output are known, and then dispatch reserves after the uncertainty is realized. Similarly, a supply-chain planner may position inventory before future demand is observed, and then route shipments or place emergency orders afterward. Two-stage stochastic programming has been studied for several decades and has developed a rich mathematical theory \citep{placeholder_stochastic_programming_textbook,placeholder_two_stage_sp_survey}. Nevertheless, it remains an active research area because many important applications still require solving large and difficult stochastic optimization problems \citep{placeholder_modern_sp_applications}.

In many modern applications, the uncertainty distribution is not fixed and known. Instead, the decision maker observes contextual information before making the first-stage decision. In the energy example, this context may include weather forecasts, historical demand patterns, and market information. In logistics, it may include calendar effects, regional conditions, or recent order data. This gives rise to contextual stochastic programming: given historical context--scenario pairs, the goal is to learn a policy that maps each context to a high-quality first-stage decision. Contextual stochastic programming therefore creates a natural opportunity for machine learning to address one of the major bottlenecks in stochastic programming: learning how uncertainty should affect decisions.

Broadly, existing approaches to contextual stochastic programming fall into three categories.

\begin{itemize}
    \item \textbf{Predict-then-optimize.} A model first learns the conditional distribution of the uncertainty, or a sample approximation of it, and a stochastic program is then solved using the learned distribution. This approach can be combined with classical optimization methods and, when the learned distribution is accurate enough, can recover an exact optimal policy for the estimated problem. Its weakness is that learning full conditional distributions can be data intensive, and small statistical errors in the learned distribution can translate into large differences in downstream decision quality \citep{placeholder_predict_then_optimize_sp,placeholder_statistical_error_decision_quality}.

    \item \textbf{Direct policy learning.} A model directly learns a map from context to first-stage decisions, bypassing explicit distribution learning. This can be effective when the optimal decision policy is simpler than the conditional distribution itself. However, direct policy learning must still ensure that the produced decisions are feasible. This is difficult for problems with complicated constraints, and especially difficult when the first-stage decision includes integer variables \citep{placeholder_direct_policy_learning,placeholder_feasible_neural_policies,placeholder_integer_constraints_learning}.

    \item \textbf{Decision-focused learning.} A model learns scenarios, costs, or other optimization inputs specifically to improve downstream decision quality. The learned input is passed to a surrogate optimization problem, and the resulting decision is evaluated on the true downstream objective. This approach naturally respects the constraints of the surrogate optimization problem. It also extends naturally to mixed-integer settings by training on continuous relaxations and using integer restrictions at inference. Its main weakness is computational: training requires repeatedly solving optimization problems and differentiating through their solution maps, which can quickly become a bottleneck \citep{placeholder_decision_focused_learning_overview,placeholder_dfl_stochastic_programming}.
\end{itemize}

This paper focuses on the third category. In particular, we study decision-focused scenario learning for contextual two-stage stochastic linear programming. A scenario generator maps each context to a synthetic scenario. The first-stage decision is then obtained by solving the corresponding single-scenario optimization problem, and the generator is trained to minimize downstream decision loss.

Learning a single scenario is attractive for two reasons. First, inference is fast: solving a single-scenario problem is much cheaper than solving a large sample-average approximation, whose deterministic equivalent grows with the number of scenarios. Second, training is also faster, since the surrogate problem must be re-solved at every training step. However, single-scenario learning raises two fundamental questions. The first is expressivity: can a single generated scenario represent the optimal policy of the full contextual stochastic program? The second is trainability: can such scenarios be learned with gradient-based methods, despite the non-differentiability of linear programming solution maps?

We address both questions. Our starting point is the observation that synthetic scenarios should be understood as decision parametrizations, not necessarily as realistic forecasts. A decision-optimal scenario may lie outside the support of the true uncertainty distribution, and may not correspond to any physically meaningful realization. This is not a defect. It is precisely what allows a low-dimensional scenario representation to encode the effect of a much richer conditional distribution.

\textbf{Our contribution.}
We make the following three contributions.

\begin{itemize}
    \item \textbf{A realizability theory for decision-focused scenario learning.}
    We prove that, under weak well-posedness assumptions, optimal contextual stochastic programming policies can be represented by solving carefully chosen single-scenario problems. In particular, a scenario generator does not need to learn a full distribution, or even a full scenario. Learning one synthetic second-stage component is often sufficient to induce the optimal first-stage decision. This gives a general theoretical explanation for why single-scenario decision-focused methods can be expressive enough for contextual stochastic programming.

    \item \textbf{Equivalence between scenario components and direct decision learning.}
    We show that learning different second-stage components can be theoretically equivalent from the perspective of induced first-stage decisions. In particular, learning costs, right-hand sides, or suitably transformed recourse matrices can encode the same optimal policy. This substantially enlarges the applicability of existing methods, because the learnable component can be chosen based on computational convenience, dimensionality, and feasibility constraints rather than on which component is random in the original problem.

    \item \textbf{Log-barrier smoothing for trainable scenario learning.}
    We study log-barrier smoothing of the single-scenario surrogate problem. The smoothing makes the scenario-to-decision map differentiable and yields closed-form implicit derivatives. More importantly, we prove that the resulting decision-focused loss has a favorable optimization geometry: it is componentwise invex, or equivalently has no spurious stationary points, with respect to the relevant learned scenario components. This provides a theoretical explanation for why log-barrier smoothing avoids the vanishing-gradient pathologies associated with other smoothing schemes. Motivated by this theory, we propose an interior-point-inspired training method that starts with a large barrier parameter and gradually decreases it during training.
\end{itemize}

Together, these results provide both an expressivity theory and an optimization theory for decision-focused single-scenario learning. The expressivity theory explains why synthetic scenarios can encode optimal policies. The smoothing theory explains why such scenarios can be learned by gradient-based methods. The resulting method combines the structure-preserving nature of mathematical programming with the flexibility of neural prediction.

\textbf{Closest approaches in the literature.}
Our work is most closely related to four strands of prior work. First, decision-focused learning has been extensively studied for single-stage optimization problems, especially problems with random cost vectors \citep{placeholder_single_stage_dfl,placeholder_spo,placeholder_decision_focused_surveys}. Some recent papers use smoothing, including log-barrier smoothing, to enable differentiation through optimization layers \citep{placeholder_log_barrier_dfl,placeholder_mandi_log_barrier}. Our setting is different because contextual two-stage stochastic programs include recourse decisions and second-stage constraints, which substantially enrich the geometry of the solution map.

Second, early work on decision-focused learning for stochastic programming trained predictive models through smoothed stochastic optimization problems \citep{placeholder_donti_stochastic_dfl}. These methods typically learn distributions or collections of scenarios and often rely on quadratic smoothing. Quadratic smoothing is convenient computationally, but it is known to suffer from vanishing-gradient behavior in regions with nonzero regret \citep{placeholder_quadratic_smoothing_vanishing_gradients}. Our use of log-barrier smoothing is motivated by the opposite behavior: the barrier preserves interior sensitivity and leads to componentwise invexity of the downstream loss.

Third, recent work has studied decision-focused scenario generation for stochastic programming \citep{placeholder_islip_scenario_generation,placeholder_scenario_generation_black_box}. These methods are close in spirit to ours, but often rely on black-box surrogate models or empirical training procedures without a general theory of when generated scenarios can represent optimal policies. Our contribution is to provide such a theory, together with explicit constructions of decision-optimal synthetic scenarios.

Fourth, several papers derive single-scenario optimality guarantees for restricted stochastic programming settings, especially when uncertainty appears only in the second-stage right-hand side or in application-specific structures \citep{placeholder_lex_spo,placeholder_homem_de_mello_single_scenario,placeholder_application_specific_single_scenario}. Our results aim to remove these restrictions. We show that decision-optimal scenario learning can be understood at the level of the induced decision map, which allows the method to handle arbitrary second-stage randomness by learning whichever synthetic component is most convenient.

The remainder of the paper is organized as follows. Section~\ref{sec:problem} introduces contextual two-stage stochastic programming and decision-focused scenario generation. Section~\ref{sec:realizability} proves the realizability and component-equivalence results. Section~\ref{sec:barrier} develops the log-barrier smoothing theory and the no-spurious-stationary-points result. Section~\ref{sec:method} presents the interior-point-inspired training method. Section~\ref{sec:experiments} evaluates the method on standard benchmarks.
